% GENERATED FILE. Edit canonical lessons, then run npm run generate:books.
% canonical-source-hash: fnv1a64:4a3b298605c66185
% canonical-lessons: TA-C141-tivu, TA-C141-pucci, TA-C141-cilanti, TA-C141-panri, TA-C141-kalutai

\chapter{Island, Insect, Spider, Pig, Donkey}
\label{ch:ta-island-insect-spider-pig-donkey}
\hlchaptermodality{141}

\begin{chapteropening}
\textbf{By the end of this chapter:} \emph{I can say an island, an insect, a spider, a pig and a donkey.}

Complete the last lesson of chapter 141: i can say an island, an insect, a spider, a pig and a donkey.
\end{chapteropening}

\section[tīvu]{\ta{தீவு} (tīvu) — an island}

\label{lesson:TA-C141-tivu}



\begin{quote}
Before the new one: say the Tamil for a waterfall, then the Tamil for a hill.
\end{quote}

\subsection*{You'll want to know: \ta{தீவு}}
\textbf{\ta{தீவு}} — \emph{tīvu} — \textquotedblleft{}an island\textquotedblright{}.

\begin{grammarlens}[title={What you've built}]
One more word for this chapter.
\end{grammarlens}

\subsection*{Your turn}
\begin{itemize}
\raggedright
  \item \emph{Say it:} \emph{tīvu}
  \item \emph{Say it:} \emph{tīvu}, once more
  \item \emph{Say it:} say \emph{tīvu}
  \item \emph{Recall:} say the Tamil for a waterfall, then the Tamil for a hill, then say \emph{tīvu} again
\end{itemize}

\begin{tcolorbox}[breakable,colback=teal!4,colframe=teal!35!black,title={Before you move on}]
What does \ta{தீவு} mean? (\textquotedblleft{}An island\textquotedblright{}.) Say it once more.
\end{tcolorbox}
\section[pūcci]{\ta{பூச்சி} (pūcci) — an insect}

\label{lesson:TA-C141-pucci}



\begin{quote}
Before the new one: say the Tamil for a hill, then the Tamil for an island.
\end{quote}

\subsection*{You'll want to know: \ta{பூச்சி}}
\textbf{\ta{பூச்சி}} — \emph{pūcci} — \textquotedblleft{}an insect\textquotedblright{}.

\begin{grammarlens}[title={What you've built}]
One more word for this chapter.
\end{grammarlens}

\subsection*{Your turn}
\begin{itemize}
\raggedright
  \item \emph{Say it:} \emph{pūcci}
  \item \emph{Say it:} \emph{pūcci}, once more
  \item \emph{Say it:} say \emph{pūcci}
  \item \emph{Recall:} say the Tamil for a hill, then the Tamil for an island, then say \emph{pūcci} again
\end{itemize}

\begin{tcolorbox}[breakable,colback=teal!4,colframe=teal!35!black,title={Before you move on}]
What does \ta{பூச்சி} mean? (\textquotedblleft{}An insect\textquotedblright{}.) Say it once more.
\end{tcolorbox}
\section[cilanti]{\ta{சிலந்தி} (cilanti) — a spider}

\label{lesson:TA-C141-cilanti}



\begin{quote}
Before the new one: say the Tamil for an island, then the Tamil for an insect.
\end{quote}

\subsection*{You'll want to know: \ta{சிலந்தி}}
\textbf{\ta{சிலந்தி}} — \emph{cilanti} — \textquotedblleft{}a spider\textquotedblright{}.

\begin{grammarlens}[title={What you've built}]
One more word for this chapter.
\end{grammarlens}

\subsection*{Your turn}
\begin{itemize}
\raggedright
  \item \emph{Say it:} \emph{cilanti}
  \item \emph{Say it:} \emph{cilanti}, once more
  \item \emph{Say it:} say \emph{cilanti}
  \item \emph{Recall:} say the Tamil for an island, then the Tamil for an insect, then say \emph{cilanti} again
\end{itemize}

\begin{tcolorbox}[breakable,colback=teal!4,colframe=teal!35!black,title={Before you move on}]
What does \ta{சிலந்தி} mean? (\textquotedblleft{}A spider\textquotedblright{}.) Say it once more.
\end{tcolorbox}
\section[paṉṟi]{\ta{பன்றி} (paṉṟi) — a pig}

\label{lesson:TA-C141-panri}



\begin{quote}
Before the new one: say the Tamil for an insect, then the Tamil for a spider.
\end{quote}

\subsection*{You'll want to know: \ta{பன்றி}}
\textbf{\ta{பன்றி}} — \emph{paṉṟi} — \textquotedblleft{}a pig\textquotedblright{}.

\begin{grammarlens}[title={What you've built}]
One more word for this chapter.
\end{grammarlens}

\subsection*{Your turn}
\begin{itemize}
\raggedright
  \item \emph{Say it:} \emph{paṉṟi}
  \item \emph{Say it:} \emph{paṉṟi}, once more
  \item \emph{Say it:} say \emph{paṉṟi}
  \item \emph{Recall:} say the Tamil for an insect, then the Tamil for a spider, then say \emph{paṉṟi} again
\end{itemize}

\begin{tcolorbox}[breakable,colback=teal!4,colframe=teal!35!black,title={Before you move on}]
What does \ta{பன்றி} mean? (\textquotedblleft{}A pig\textquotedblright{}.) Say it once more.
\end{tcolorbox}
\section[kaḻutai]{\ta{கழுதை} (kaḻutai) — a donkey}

\label{lesson:TA-C141-kalutai}



\begin{quote}
Before the new one: say the Tamil for a spider, then the Tamil for a pig.
\end{quote}

\subsection*{You'll want to know: \ta{கழுதை}}
\textbf{\ta{கழுதை}} — \emph{kaḻutai} — \textquotedblleft{}a donkey\textquotedblright{}.

\begin{grammarlens}[title={What you've built}]
One more word for this chapter.
\end{grammarlens}

\subsection*{Your turn}
\begin{itemize}
\raggedright
  \item \emph{Say it:} \emph{kaḻutai}
  \item \emph{Say it:} \emph{kaḻutai}, once more
  \item \emph{Say it:} say \emph{kaḻutai}
  \item \emph{Recall:} say the Tamil for a spider, then the Tamil for a pig, then say \emph{kaḻutai} again
\end{itemize}

\begin{tcolorbox}[breakable,colback=teal!4,colframe=teal!35!black,title={Before you move on}]
What does \ta{கழுதை} mean? (\textquotedblleft{}A donkey\textquotedblright{}.) Say it once more.
\end{tcolorbox}
